\section{Conclusion and Future Work}

In this work, we successfully surpassed the limitations posed by traditional single-GPU-per-model training by implementing model parallelism with only a few modifications to the existing PyTorch transformer implementations. We efficiently trained transformer based models up to 8.3 billion parameter on 512 NVIDIA V100 GPUs with 8-way model parallelism and achieved up to 15.1 PetaFLOPs sustained over the entire application. We also showed that for BERT models, careful attention to the placement of layer normalization in BERT-like models is critical to achieving increased accuracies as the model size increases. We study the effect of model size on down-stream task accuracy and achieve far superior results on downstream tasks and establish new SOTA for WikiText103, LAMBADA, and RACE datasets. Finally, we open sourced our code to enable future work leveraging model parallel transformers.

There are several directions for future work. Continuing to increase the scale of pretraining is a promising line of investigation that will further test existing deep learning hardware and software. To realize this, improvements in the efficiency and memory footprint of optimizers will be needed. In addition, training a model with more than 16 billion parameters will demand more memory than is available within 16 GPUs of a DGX-2H box. For such models, a hybrid intra-layer and inter-layer model parallelism along with inter-node model parallelism would be more suitable.
%Increasing the scale of pretraining and transfer is not the only way to demonstrate the effectiveness of large scale language modeling. To this end 
Three other directions of investigation include (a) pretraining different model families (XLNet, T5), (b) evaluating performance of large models across more difficult and diverse downstream tasks (e.g. Generative Question Answering, Summarization, and Conversation), and (c) using knowledge distillation to train small student models from these large pretrained teacher models.
